% Ροή του αλγορίθμου δύο φάσεων του Kociemba (περιεχόμενο σχήματος· το
% περιβάλλον figure και η λεζάντα ορίζονται στο κεφάλαιο 4). Φασικά βάθη
% κατά τα κλειδωμένα δεδομένα: φάση 1 <= 12, φάση 2 <= 18.
\begin{tikzpicture}[
  state/.style={draw, rectangle, align=center, font=\small,
                minimum height=1.1cm, minimum width=2.7cm, inner sep=5pt},
  phase/.style={font=\footnotesize},
  desc/.style={font=\scriptsize, align=center, anchor=north},
  flow/.style={-{Stealth[length=2.8mm]}, semithick},
  node distance=2.6cm
]
\node[state] (start) {Αυθαίρετη\\ κατάσταση};
\node[state, right=of start] (g1) {Υποομάδα \(G_1\)};
\node[state, right=of g1] (solved) {Λυμένη\\ κατάσταση};

\draw[flow] (start) -- node[phase, above] {Φάση 1} (g1);
\draw[flow] (g1) -- node[phase, above] {Φάση 2} (solved);

\node[desc] at ($(start.east)!0.5!(g1.west) + (0,-0.7)$)
  {IDA* στις συντεταγμένες\\ φάσης 1, βάθος \(\le 12\)};
\node[desc] at ($(g1.east)!0.5!(solved.west) + (0,-0.7)$)
  {IDA* στις συντεταγμένες\\ φάσης 2, βάθος \(\le 18\)};

\draw[flow, dashed] (solved.north) -- ++(0,0.85) -| (start.north);
\node[font=\scriptsize, align=center, anchor=south]
  at ($(start.north)!0.5!(solved.north) + (0,0.95)$)
  {επαναληπτική εμβάθυνση στο μήκος της φάσης 1\\
   \(\rightarrow\) διαδοχικά καλύτερες συνολικές λύσεις};
\end{tikzpicture}
